% =====================================================================
% AegisBench, WACV 2027 Evaluations and Datasets track
% BODY ONLY. Use the official WACV 2027 author kit for the preamble,
% \documentclass, and the wacv.sty it ships. Do not substitute a
% hand-written preamble: a paper not using the template is rejected
% without review.
%
% Every number in the tables below is measured and traceable to a file
% in results/. Do not edit a number here without changing its source.
% \TODO markers show what is still to be written.
% =====================================================================

\newcommand{\TODO}[1]{\textcolor{red}{[TODO: #1]}}

% ---------------------------------------------------------------------
\section{Introduction}
% ---------------------------------------------------------------------
\TODO{Paragraph 1, stakes. Aerial SAR, why automated detection matters,
why failure is costly. Two sentences of context then technical framing.}

\TODO{Paragraph 2, the gap. Published results on HERIDAL and SARD are
strong; both were collected in calm clear conditions; deployments are
not. State plainly that no systematic measurement exists of behaviour
under the visual conditions of the disasters these systems respond to.
This is the sentence the paper hangs on.}

\TODO{Paragraph 3, the approach. Three things that make this more than
an ImageNet-C transplant: corruptions modeled on documented disaster
optics, severities calibrated against measurable statistics, frozen
clean-validation operating points.}

\TODO{Paragraph 4, the findings. Lead with the low-light collapse and
its interval. Then the spread across conditions. Then the localization
decoupling. One sentence each.}

% Paragraph 5, contributions: ALREADY WRITTEN. Paste here.

% ---------------------------------------------------------------------
\section{Related Work}
% ---------------------------------------------------------------------

\subsection{Aerial and Search-and-Rescue Person Detection}
\TODO{NOT WRITTEN. Highest priority prose remaining. Must cite HERIDAL
and SARD (the datasets you evaluate on), plus TinyPerson and
SeaDronesSee, both WACV papers in adjacent SAR domains. Must contain the
SARD-Corr differentiation paragraph: see docs/WRITING_PACK.md section 3
for text. Ends by establishing that all of this is evaluated under clear
conditions.}

\subsection{Corruption Robustness Benchmarking}
% ALREADY WRITTEN. Paste here.

\subsection{Adverse Weather and Low-Light Vision}
% ALREADY WRITTEN. Paste here. Candidate for the largest cut, see
% docs/DRAFT_REVIEW.md section 5.

\subsection{Small-Object Detection and Tiling}
% ALREADY WRITTEN. Paste here. Trim the FPN background paragraph.

% ---------------------------------------------------------------------
\section{Benchmark Design}
% ---------------------------------------------------------------------
% ALREADY WRITTEN: taxonomy, physical models, calibration, determinism,
% ground-truth alignment. Paste here, WITH the Koschmieder correction
% from docs/DRAFT_REVIEW.md section 1.1 applied.

\begin{table*}[t]
\centering\small
\caption{The nine corruptions, their disaster family, the optical
mechanism each models, and the statistic defining its severity ladder.
Arrows give the direction the statistic must move as severity increases.
Values are means over a fixed random subsample of 20 HERIDAL test
images, held constant across all corruptions and severities;
monotonicity is verified automatically on these means.
$\dagger$~\texttt{rain\_streaks} has no closed-form image statistic
isolating streak density, so this row reports the constructed generation
parameter rather than a measured image statistic; see
Section~\ref{sec:calibration}.}
\label{tab:taxonomy}
\begin{tabular}{llllrrr}
\toprule
Corruption & Family & Optical mechanism & Statistic & s1 & s2 & s3 \\
\midrule
\texttt{water\_glare}     & Flood      & Specular sun-glint saturation         & saturated frac.\ $\uparrow$      & 0.0116 & 0.0182 & 0.1895 \\
\texttt{turbidity\_cast}  & Flood      & Sediment-laden water, contrast loss   & RMS contrast $\downarrow$        & 0.1098 & 0.0701 & 0.0378 \\
\texttt{inundation}       & Flood      & Semi-transparent standing water       & occluded frac.\ $\uparrow$       & 0.0830 & 0.2235 & 0.3921 \\
\midrule
\texttt{smoke\_haze}      & Wildfire   & Koschmieder, neutral-gray airlight    & RMS contrast $\downarrow$        & 0.1170 & 0.0776 & 0.0475 \\
\texttt{fire\_warm\_tint} & Wildfire   & Low-CCT firelight, warm airlight haze & red/blue ratio $\uparrow$        & 1.6743 & 2.0721 & 2.4516 \\
\midrule
\texttt{rain\_streaks}    & Storm      & Directional streaks over a rain veil  & streaks/MP $\uparrow$~$\dagger$  & 120    & 350    & 800    \\
\texttt{motion\_blur}     & Storm      & Wind-induced platform shake           & edge strength $\downarrow$       & 0.0764 & 0.0572 & 0.0472 \\
\texttt{low\_light}       & Storm      & Linear-domain photon scaling, noise   & mean luminance $\downarrow$      & 0.2355 & 0.1533 & 0.1175 \\
\midrule
\texttt{dust\_haze}       & Earthquake & Koschmieder, brown mineral airlight   & RMS contrast $\downarrow$        & 0.1094 & 0.0720 & 0.0489 \\
\bottomrule
\end{tabular}
\end{table*}

\begin{figure*}[t]
\centering
\includegraphics[width=\textwidth]{figures/taxonomy_s3.pdf}
\caption{The nine corruptions at severity 3 on a cropped HERIDAL frame.
Each models a documented optical signature of a disaster environment
rather than a generic distortion. Severity is calibrated per corruption
against a declared image statistic (Table~\ref{tab:taxonomy}).}
\label{fig:taxonomy}
\end{figure*}

\subsection{Calibration and Severity}
\label{sec:calibration}
% ALREADY WRITTEN and already corrected for the sampling procedure.

% ---------------------------------------------------------------------
\section{Evaluation Protocol}
% ---------------------------------------------------------------------
% ALREADY WRITTEN: frozen operating points, shared evaluator, tiling,
% bootstrap, provenance. Paste here. NOTE: delete the duplicated
% sentence in the tiling subsection (the "removes tiling itself as a
% confound" sentence appears twice).

% ---------------------------------------------------------------------
\section{Experimental Setup}
% ---------------------------------------------------------------------
% ALREADY DRAFTED with bracketed placeholders. Fill them from the tables
% below and from docs/WRITING_PACK.md section 1.

\begin{table}[t]
\centering\small
\caption{Training configuration. Hyperparameters follow established
per-architecture fine-tuning practice rather than a shared schedule.
Corrupted data enters neither training nor threshold selection.}
\label{tab:training}
\begin{tabular}{lccc}
\toprule
 & YOLOv11 & RT-DETR & Faster R-CNN \\
\midrule
Base weights   & \texttt{yolo11m} & \texttt{rtdetr-l} & torchvision \\
Epochs         & 80     & 80      & 26     \\
Batch size     & 4      & 2       & 2      \\
Input size     & 1024   & 1024    & 1024   \\
Initial LR     & 0.01   & 0.0001  & 0.005  \\
Patience       & 20     & 20      & ---    \\
Seed           & 17     & 17      & 17     \\
\bottomrule
\end{tabular}
\end{table}

\begin{table}[t]
\centering\small
\caption{Frozen operating points, selected once per model per dataset by
maximizing F1 on clean validation data and held fixed across all 28
conditions for that model.}
\label{tab:thresholds}
\begin{tabular}{lcc}
\toprule
Model & HERIDAL & SARD \\
\midrule
Faster R-CNN & 0.99 & 0.96 \\
RT-DETR      & 0.65 & 0.60 \\
YOLOv11      & 0.45 & 0.30 \\
\bottomrule
\end{tabular}
\end{table}

% ---------------------------------------------------------------------
\section{Results}
% ---------------------------------------------------------------------
% NOT WRITTEN. All numbers below are measured. Prose guidance for each
% subsection is in docs/WRITING_PACK.md section 2.

\subsection{Clean Baselines}
\TODO{Prose. Establishes the pipeline is sound before any robustness
claim. Compare against the published HERIDAL reference point and note
that modest deltas across architectures are expected. Note SARD is the
easier task, which is why everything downstream is reported as relative
drop.}

\begin{table}[t]
\centering\small
\caption{Clean-condition performance. All degradation results are
reported relative to these baselines, since the two datasets do not
share a difficulty floor.}
\label{tab:clean}
\begin{tabular}{llcccc}
\toprule
Model & Dataset & R & P & mAP$_{50}$ & mAP$_{50:95}$ \\
\midrule
Faster R-CNN & HERIDAL & 0.745 & 0.903 & 0.785 & 0.518 \\
RT-DETR      & HERIDAL & 0.792 & 0.881 & 0.835 & 0.556 \\
YOLOv11      & HERIDAL & 0.769 & 0.855 & 0.803 & 0.510 \\
\midrule
Faster R-CNN & SARD    & 0.879 & 0.967 & 0.896 & 0.615 \\
RT-DETR      & SARD    & 0.888 & 0.972 & 0.929 & 0.633 \\
YOLOv11      & SARD    & 0.891 & 0.955 & 0.928 & 0.601 \\
\bottomrule
\end{tabular}
\end{table}

\subsection{Robustness Across Conditions}
\TODO{Prose. The spread is the story: rain costs 3.6 points at severity
3 while low light costs everything. A single averaged robustness score
would hide the entire finding.}

\begin{figure*}[t]
\centering
\includegraphics[width=\textwidth]{figures/heatmap_sard.pdf}
\caption{Relative recall drop, $(\text{clean}-\text{corrupted})/\text{clean}$,
for each detector on SARD. Rows are grouped by disaster family and
ordered to match Table~\ref{tab:taxonomy}. Darker cells indicate greater
degradation. The \texttt{low\_light} row saturates for all three
architectures.}
\label{fig:heatmap}
\end{figure*}

\begin{figure*}[t]
\centering
\includegraphics[width=\textwidth]{figures/severity_sard_yolo11.pdf}
\caption{Recall against corruption severity for YOLOv11 on SARD, with
bootstrap 95\% confidence intervals shaded (1{,}000 resamples per
condition). All curves originate from the same clean baseline. Colour
indicates disaster family; line style distinguishes corruptions within a
family.}
\label{fig:severity}
\end{figure*}

\begin{table}[t]
\centering\small
\caption{Recall by corruption and severity, YOLOv11 on SARD, bootstrap
means over 1{,}000 resamples. Clean baseline 0.891. Ordered by severity-3
recall.}
\label{tab:spectrum}
\begin{tabular}{llccc}
\toprule
Corruption & Family & s1 & s2 & s3 \\
\midrule
\texttt{rain\_streaks}    & Storm      & 0.893 & 0.883 & \textbf{0.855} \\
\texttt{water\_glare}     & Flood      & 0.859 & 0.809 & 0.666 \\
\texttt{inundation}       & Flood      & 0.860 & 0.744 & 0.494 \\
\texttt{motion\_blur}     & Storm      & 0.792 & 0.613 & 0.477 \\
\texttt{smoke\_haze}      & Wildfire   & 0.810 & 0.667 & 0.391 \\
\texttt{turbidity\_cast}  & Flood      & 0.796 & 0.576 & 0.191 \\
\texttt{fire\_warm\_tint} & Wildfire   & 0.689 & 0.311 & 0.099 \\
\texttt{dust\_haze}       & Earthquake & 0.673 & 0.272 & 0.058 \\
\texttt{low\_light}       & Storm      & 0.184 & 0.007 & \textbf{0.000} \\
\bottomrule
\end{tabular}
\end{table}

\begin{table}[t]
\centering\small
\caption{Mean relative recall drop by disaster family, averaged over
severities. Flood is the mildest family on both datasets for every
model; storm is worst on HERIDAL and earthquake on SARD.}
\label{tab:family}
\begin{tabular}{lccc}
\toprule
Family & Faster R-CNN & RT-DETR & YOLOv11 \\
\midrule
\multicolumn{4}{l}{\emph{HERIDAL}} \\
Storm      & 0.784 & 0.874 & 0.829 \\
Earthquake & 0.534 & 0.568 & 0.571 \\
Wildfire   & 0.335 & 0.459 & 0.405 \\
Flood      & 0.301 & 0.231 & 0.246 \\
\midrule
\multicolumn{4}{l}{\emph{SARD}} \\
Earthquake & 0.481 & 0.630 & 0.624 \\
Storm      & 0.401 & 0.440 & 0.413 \\
Wildfire   & 0.327 & 0.451 & 0.445 \\
Flood      & 0.264 & 0.293 & 0.252 \\
\bottomrule
\end{tabular}
\end{table}

\subsection{Detection Collapse Under Low Light}
\TODO{Prose. The headline. Give it room. Cross-dataset,
cross-architecture, with the threshold-independent mAP confirmation.
State the two safe claims and avoid the unsafe one: Faster R-CNN is
consistently most robust; the YOLOv11 versus RT-DETR ordering is
dataset-dependent and must not be generalized. Note that Faster R-CNN
has the strictest threshold and is nonetheless most robust, which closes
off the obvious alternative explanation.}

\begin{table}[t]
\centering\small
\caption{Recall under \texttt{low\_light} at each model's frozen
operating point, with bootstrap 95\% intervals where computed. Recall
reaches exactly zero for all three architectures on both datasets at
severity 3, and every one of 1{,}000 resamples agrees, so the interval is
degenerate at zero.}
\label{tab:lowlight}
\begin{tabular}{llccc}
\toprule
Model & Thr. & s1 & s2 & s3 \\
\midrule
\multicolumn{5}{l}{\emph{HERIDAL} ($n=101$)} \\
RT-DETR      & 0.65 & 0.000 [0.000, 0.000] & 0.000 [0.000, 0.000] & 0.000 [0.000, 0.000] \\
YOLOv11      & 0.45 & 0.048 [0.027, 0.073] & 0.000 [0.000, 0.000] & 0.000 [0.000, 0.000] \\
Faster R-CNN & 0.99 & 0.452 [0.365, 0.540] & 0.110 [0.064, 0.163] & 0.000 [0.000, 0.000] \\
\midrule
\multicolumn{5}{l}{\emph{SARD} ($n=862$)} \\
Faster R-CNN & 0.96 & 0.419 & 0.077 & 0.000 \\
RT-DETR      & 0.60 & 0.260 & 0.013 & 0.000 \\
YOLOv11      & 0.30 & 0.184 [0.157, 0.212] & 0.007 [0.003, 0.012] & 0.000 [0.000, 0.000] \\
\bottomrule
\end{tabular}
\end{table}

\begin{table}[t]
\centering\small
\caption{Threshold-independent confirmation on HERIDAL. mAP$_{50}$ scans
all confidence thresholds, so it is immune to the objection that a zero
reflects a badly placed operating point. It collapses in lockstep with
recall: the detections are absent, not merely below threshold.}
\label{tab:calibration}
\begin{tabular}{lcc}
\toprule
Model & Clean mAP$_{50}$ & \texttt{low\_light} s1 mAP$_{50}$ \\
\midrule
RT-DETR      & 0.835 & \textbf{0.001} \\
YOLOv11      & 0.803 & 0.064 \\
Faster R-CNN & 0.785 & 0.537 \\
\bottomrule
\end{tabular}
\end{table}

\subsection{Localization Stability}
\TODO{Prose. State the finding as a claim about magnitude, not
independence: both recall and stability decline, but recall spans nearly
the full unit range while stability never falls below 0.78. n\_common
tracks recall exactly. The operational implication: an operator can
trust the boxes that appear; the risk is the ones that never do.}

\begin{figure}[t]
\centering
\includegraphics[width=\columnwidth]{figures/decoupling_sard_yolo11.pdf}
\caption{Recall against localization stability, YOLOv11 on SARD, one
point per corruption and severity, point area proportional to the number
of instances detected in both the clean and corrupted conditions. Both
axes span the full unit range. Points sweep the entire recall range while
remaining in a narrow band near the top of the stability axis.}
\label{fig:decoupling}
\end{figure}

\begin{figure}[t]
\centering
\includegraphics[width=\columnwidth]{figures/gallery_lowlight.pdf}
\caption{\TODO{Qualitative pair. Clean frame beside the same frame under
\texttt{low\_light}, ground truth in cyan, retained detections in green.}}
\label{fig:qualitative}
\end{figure}

% ---------------------------------------------------------------------
\section{Limitations}
% ---------------------------------------------------------------------
\TODO{Seven items, full prose in docs/WRITING_PACK.md section 4:
synthetic corruptions on real imagery; single training seed; scope of
two datasets and three detectors; the rain\_streaks proxy confound;
monotonicity verified on a subsample by mean; the inundation ripple; no
mitigation study.}

% ---------------------------------------------------------------------
\section{Conclusion}
% ---------------------------------------------------------------------
\TODO{Three moves, no new numbers. Restate the gap. Restate the finding
at its strongest. State the implication and the next step. Do not close
on a claim of generality the study does not support.}

{\small
\bibliographystyle{ieee_fullname}
\bibliography{aegisbench}
}
